% Append this section before \end{document} in the master AION document.
\section{Signed Agent and Capability Packages}
\label{sec:capability-packages}

The capability-package layer turns specialist prompts, procedures and tools into portable,
inspectable additions to a customer-owned AION brain.  A package may add a sales adviser,
finance procedure, industry workflow or another bounded skill, but it cannot contain or
replace the customer's brain, Business Map or private data.

\subsection{Package Identity and Supply Chain}

Every release has an immutable package identifier and version.  Its manifest binds the
artefact's SHA--256 hash and byte length to a trusted Ed25519 publisher, package name,
department, industry, permissions, limits, tests and publishing terms.  Installation verifies
the publisher, canonical manifest hash, signature and real artefact bytes.  Reusing a version
with changed content is rejected.  Manifests declaring copied customer brain or customer data
are rejected before installation.

\subsection{Declared Least Authority}

Permissions are explicit allow-lists for network hosts, file roots, identity scopes, tools
and named execution actions.  The authorization broker begins with no authority and issues a
hashed grant only when the package is active, the exact target is declared and its cumulative
authorization ceiling has not been reached.  File paths are resolved before comparison so that
relative traversal cannot escape an allowed root.  Raw commands cannot be substituted for a
named execution action.  The package never obtains authority merely because an LLM proposed it.

The broker returns permission to a separate governed executor; it does not run arbitrary
third-party code in the mother process.  A production operating-system or container sandbox is
still required before arbitrary marketplace code can execute.  This distinction prevents a
logical allow-list from being misrepresented as process isolation.

\subsection{Qualification and Activation}

Three independent test records are mandatory:
\begin{enumerate}
  \item a dry run demonstrating expected inputs and outputs without external mutation;
  \item adversarial permission tests proving undeclared access is denied; and
  \item outcome tests establishing useful task behaviour.
\end{enumerate}

Any failed test blocks the release.  Activation is unavailable until all three pass and a
named human owner supplies an approval receipt.  The receipt is retained as a hash rather than
raw approval content.  Every runtime grant identifies the accountable owner, making a package
an owned delegation rather than an autonomous legal or managerial actor.

\subsection{Independent Component Versions}

Prompt, procedure, tool and policy components have separate declared version lists and active
selections.  Changing or rolling back one component does not require changing the underlying
model or other components.  Each selection preserves the previous version and timestamp in
component history.  This makes it possible to determine whether a behaviour change came from
a model, prompt, operating procedure, tool implementation or policy.

\subsection{Catalogue and Commercial Controls}

The catalogue distinguishes Tessaris-curated, customer-created and third-party publishers.
It presents package status, department, industry, curation state, commercial terms and the
accountable owner while explicitly excluding customer brain content.  Third-party releases
must declare commercial terms, creating a foundation for later publisher onboarding and
revenue sharing without allowing payment status to bypass safety qualification.

\subsection{Revocation, Removal and Brain Integrity}

Revocation records a reason and immediately ends authorization.  Removal changes the package
state to removed while retaining qualification, activation and revocation evidence.  Neither
operation deletes or rewrites external business state.  Automated tests preserve a control
Business Map file before and after the full activate--revoke--remove sequence to verify this
separation.

\subsection{Implemented and Remaining Work}

Pilot Fabric~0.61.0 and AION Sovereign Brain Bootstrap~0.9.0 implement trusted manifests,
artefact verification, immutable releases, mandatory qualification, human-accountable
activation, least-authority grants, bounded use, independent component selection, catalogue
terms, revocation and evidence-preserving removal.

Two Phase~6 items remain open.  First, arbitrary third-party code requires qualified native
or container sandboxes for every supported operating environment; the current authority broker
does not execute such code.  Second, actual signed department and industry pack libraries must
be authored and tested without copying customer memory.  Marketplace payment, publisher identity
review and dispute handling also remain commercial-readiness work even though publishing class
and revenue terms are already represented in the catalogue.
